% Kiểu dáng chung của sổ tay, dùng cho mọi phiên bản.
% (The handbook's shared styles, for every version.)

% --- Tiếng Việt với pdfLaTeX
\usepackage[T5]{fontenc}
\usepackage[utf8]{inputenc}
\usepackage[vietnamese]{babel}

% --- Phông chữ: Charter cho thân bài, Heros (giống Helvetica) cho tiêu đề và
% tên nút. Cả hai đều có đủ dấu tiếng Việt.
\renewcommand{\rmdefault}{bch}
\renewcommand{\sfdefault}{qhv}
\linespread{1.08}

% Ký hiệu của giao diện, gõ thẳng trong bài (✎ ⠿ ⋯ ▶ ⚠ ✓, mũi tên, chữ Hy Lạp)
\usepackage{amssymb}
\usepackage{pifont}
\DeclareUnicodeCharacter{270E}{\ding{46}}
\DeclareUnicodeCharacter{2713}{\ding{51}}
\DeclareUnicodeCharacter{283F}{\raisebox{-0.2ex}{$\vdots\!\vdots$}}
\DeclareUnicodeCharacter{22EF}{\ensuremath{\cdots}}
\DeclareUnicodeCharacter{25B6}{\ensuremath{\blacktriangleright}}
\DeclareUnicodeCharacter{26A0}{\ensuremath{\blacktriangle}}
\DeclareUnicodeCharacter{211D}{\ensuremath{\mathbb{R}}}
\DeclareUnicodeCharacter{03BB}{\ensuremath{\lambda}}
\DeclareUnicodeCharacter{03A9}{\ensuremath{\Omega}}
\DeclareUnicodeCharacter{03B1}{\ensuremath{\alpha}}
\DeclareUnicodeCharacter{2191}{\ensuremath{\uparrow}}
\DeclareUnicodeCharacter{2193}{\ensuremath{\downarrow}}
\DeclareUnicodeCharacter{2190}{\ensuremath{\leftarrow}}
\DeclareUnicodeCharacter{2192}{\ensuremath{\rightarrow}}

% --- Trang, màu, hình
\usepackage[a4paper,top=2.6cm,bottom=2.6cm,left=2.5cm,right=2.5cm,
  headheight=14pt,headsep=0.9cm,footskip=1.2cm]{geometry}
\usepackage{xcolor}
\usepackage{graphicx}
\usepackage{tikz}
\usepackage[most]{tcolorbox}
\usepackage{enumitem}
\usepackage{booktabs}
\usepackage{colortbl}
\usepackage{amsmath}
\usepackage{fancyvrb}
\usepackage{fancyhdr}
\usepackage[hidelinks]{hyperref}

% Bảng màu theo giao diện SumiTeX
\definecolor{muc}{HTML}{1F2430}       % mực: chữ chính
\definecolor{sumi}{HTML}{3B5BDB}      % màu nhấn
\definecolor{sumi-dam}{HTML}{2B44B0}
\definecolor{sumi-nhat}{HTML}{EEF2FF}
\definecolor{meo}{HTML}{2F9E44}
\definecolor{meo-nhat}{HTML}{EFFAF1}
\definecolor{chuy}{HTML}{D9480F}
\definecolor{chuy-nhat}{HTML}{FFF4E6}
\definecolor{chu-mo}{HTML}{6B7280}    % chữ phụ
\definecolor{vien}{HTML}{E3E6EC}      % đường viền nhạt
\definecolor{vien-dam}{HTML}{C5CAD3}
\definecolor{nen-ma}{HTML}{F5F6F8}    % nền khung mã
\AtBeginDocument{\color{muc}}

% --- Đoạn văn và danh sách
\setlength{\parskip}{5pt}
\setlength{\parindent}{0pt}

% Danh sách canh trái (không dàn đều) để tên nút dài không làm giãn chữ
\setlist{itemsep=3pt,topsep=4pt,leftmargin=1.6em,before=\raggedright}
\setlist[itemize]{label={\color{sumi}\small\raisebox{0.1ex}{\textbullet}}}
\setlist[enumerate]{label={\color{sumi}\sffamily\bfseries\arabic*.}}
\setlist[description]{font={\sffamily\bfseries\color{muc}},leftmargin=1.2em}
% Cho phép giãn dòng thêm một chút để đoạn mã dài không tràn lề
\setlength{\emergencystretch}{3em}
% Logo LaTeX căn lại cho hợp với phông Charter
\DeclareRobustCommand{\LaTeX}{L\kern-.22em\raisebox{.32ex}{\scriptsize A}\kern-.1em\TeX}
% Bảng: đường kẻ nhạt, hàng đầu chữ không chân
\arrayrulecolor{vien-dam}
\renewcommand{\arraystretch}{1.25}

% --- Đầu trang, chân trang
\pagestyle{fancy}
\fancyhf{}
\fancyhead[L]{\sffamily\footnotesize\color{chu-mo}Sổ tay SumiTeX}
\fancyhead[R]{\sffamily\footnotesize\color{chu-mo}\leftmark}
\fancyfoot[C]{\sffamily\footnotesize\color{chu-mo}\thepage}
\renewcommand{\headrule}{{\color{vien}\hrule height 0.6pt}}
\renewcommand{\sectionmark}[1]{\markboth{#1}{}}
\fancypagestyle{plain}{\fancyhf{}\renewcommand{\headrule}{}%
  \fancyfoot[C]{\sffamily\footnotesize\color{chu-mo}\thepage}}

% --- Tiêu đề: số chương nằm trong một ô màu nhấn, kẻ một đường mảnh bên dưới
\usepackage{titlesec}
\newcommand{\ochuong}[1]{\tikz[baseline=(s.base)]\node[fill=sumi,text=white,
  rounded corners=4pt,inner xsep=7pt,inner ysep=4pt,minimum width=2em,
  font=\sffamily\large\bfseries](s){#1};}
\titleformat{\section}[block]{\sffamily\LARGE\bfseries\color{muc}}
  {\ochuong{\thesection}}{0.55em}{}[\vspace{6pt}{\color{vien}\titlerule[0.8pt]}]
\titleformat{name=\section,numberless}[block]{\sffamily\LARGE\bfseries\color{muc}}
  {}{0pt}{}[\vspace{6pt}{\color{vien}\titlerule[0.8pt]}]
\titlespacing*{\section}{0pt}{22pt plus 4pt minus 2pt}{12pt}
\titleformat{\subsection}{\sffamily\large\bfseries\color{muc}}
  {\textcolor{sumi}{\thesubsection}}{0.6em}{}
\titlespacing*{\subsection}{0pt}{16pt plus 3pt minus 2pt}{6pt}

% --- Mục lục
\usepackage[titles]{tocloft}
\setlength{\cftbeforesecskip}{0.5em}
\renewcommand{\cftsecfont}{\sffamily\bfseries}
\renewcommand{\cftsecpagefont}{\sffamily\bfseries\color{sumi}}
\renewcommand{\cftsecleader}{\hfill}
\renewcommand{\cftsubsecfont}{\color{muc}}
\renewcommand{\cftsubsecpagefont}{\color{chu-mo}}
\renewcommand{\cftdot}{\textcolor{vien-dam}{.}}
\setlength{\cftsecnumwidth}{2em}

% --- Khung: Mẹo, Lưu ý, Thực hành
% Cùng một dáng: nền nhạt, vạch màu bên trái, nhãn nhỏ có biểu tượng ở trên.
\newcommand{\bieutuong}[2]{\tikz[baseline=-0.6ex]\node[circle,fill=#1,text=white,
  inner sep=0pt,minimum size=1.15em,font=\sffamily\bfseries\scriptsize]{#2};}
% Bút chì đặc, vẽ bằng tikz cho cùng dáng với ★ và ! (ký tự ✎ quá mảnh)
\newcommand{\bieutuongbut}[1]{\tikz[baseline=-0.6ex]{%
  \fill[#1] (0,0) circle (0.575em);
  \begin{scope}[rotate=45,white]
    \fill (-0.30em,-0.10em) rectangle (0.20em,0.10em);   % thân
    \fill (-0.42em,-0.10em) rectangle (-0.34em,0.10em);  % tẩy
    \fill (0.24em,-0.10em) -- (0.42em,0) -- (0.24em,0.10em) -- cycle; % ngòi
  \end{scope}}}
% \nhankhung{màu}{ký tự trong vòng tròn}{nhãn}; \nhankhungve{màu}{hình}{nhãn}
\newcommand{\nhankhungve}[3]{{\sffamily\bfseries\small\color{#1}#2\
  #3}\par\vspace{3pt}}
\newcommand{\nhankhung}[3]{\nhankhungve{#1}{\bieutuong{#1}{#2}}{#3}}
\tcbset{khung/.style={enhanced,breakable,frame hidden,boxrule=0pt,arc=5pt,
  colback=#1-nhat,borderline west={3pt}{0pt}{#1},
  left=12pt,right=12pt,top=8pt,bottom=9pt,before skip=12pt,after skip=12pt,
  parbox=false}}
\newtcolorbox{meo}{khung=meo,before upper={\nhankhung{meo}{\ding{72}}{Mẹo}}}
\newtcolorbox{chuy}{khung=chuy,before upper={\nhankhung{chuy}{!}{Lưu ý}}}
\newtcolorbox{thuchanh}[1][]{khung=sumi,
  before upper={\nhankhungve{sumi}{\bieutuongbut{sumi}}{Thực hành\ifx\relax#1\relax\else
    \ \textcolor{chu-mo}{·}\ #1\fi}}}

% --- Khung mã: nền xám nhạt, bo góc. Dùng \begin{ma} … \end{ma}
\newtcolorbox{khungma}{enhanced,breakable,frame hidden,boxrule=0pt,arc=5pt,
  colback=nen-ma,left=10pt,right=10pt,top=6pt,bottom=6pt,
  before skip=8pt,after skip=10pt}
\newenvironment{ma}{\VerbatimEnvironment\begin{khungma}%
  \begin{Verbatim}[fontsize=\small,formatcom=\color{muc}]}%
  {\end{Verbatim}\end{khungma}}

% --- Phím và nút
% \phim{Ctrl}+\phim{S}: vẽ một phím. \nut{Download}: tên một nút trên màn hình.
\newcommand{\phim}[1]{\tikz[baseline=(k.base)]{\node[fill=vien-dam,
  rounded corners=2.5pt,inner xsep=3.5pt,inner ysep=1.6pt,yshift=-0.7pt,
  font=\footnotesize\sffamily]{\strut\phantom{#1}};
  \node[draw=vien-dam,fill=white,rounded corners=2.5pt,inner xsep=3.5pt,
  inner ysep=1.6pt,font=\footnotesize\sffamily,text=muc](k){\strut#1};}}
\newtcbox{\nut}{on line,enhanced,frame hidden,boxrule=0pt,arc=3pt,
  colback=sumi-nhat,coltext=sumi-dam,boxsep=0pt,left=2.2pt,right=2.2pt,
  top=1.2pt,bottom=1.2pt,fontupper=\sffamily\bfseries\small}
\newcommand{\lenh}[1]{\texttt{\textcolor{sumi-dam}{\textbackslash#1}}}

% --- Ảnh chụp màn hình
% \khunganh{độ rộng}{tệp}: ảnh trong một khung bo góc, viền nhạt
\newcommand{\khunganh}[2]{\tikz\node[draw=vien-dam,line width=0.6pt,
  rounded corners=5pt,inner sep=3pt,fill=white]
  {\includegraphics[width=\dimexpr#1-6.6pt\relax]{#2}};}
\newcommand{\chuthichanh}[1]{\par\vspace{4pt}{\sffamily\small\color{chu-mo}#1}}
% \anh[độ rộng]{tên tệp}{chú thích}
\newcommand{\anh}[3][\linewidth]{%
  \begin{center}
    \khunganh{#1}{#2}\chuthichanh{#3}
  \end{center}}

% Một con số tròn đặt lên ảnh, toạ độ từ 0 đến 1 tính từ góc trái dưới:
% \chuthich{x}{y}{số}
\newcommand{\chuthich}[3]{%
  \node[circle,fill=sumi,text=white,font=\footnotesize\bfseries\sffamily,
    inner sep=1.5pt,minimum size=14pt,draw=white,line width=1pt]
    at (#1,#2) {#3};}

% \anhcochuthich{độ rộng}{tệp}{các \chuthich}{chú thích}
\newcommand{\anhcochuthich}[4]{%
  \begin{center}
    \begin{tikzpicture}
      \node[anchor=south west,inner sep=0] (anh) at (0,0)
        {\includegraphics[width=\dimexpr#1-10pt\relax]{#2}};
      \draw[vien-dam,line width=0.6pt,rounded corners=5pt]
        ([shift={(-3pt,-3pt)}]anh.south west) rectangle
        ([shift={(3pt,3pt)}]anh.north east);
      \begin{scope}[x={(anh.south east)},y={(anh.north west)}]
        #3
      \end{scope}
    \end{tikzpicture}\chuthichanh{#4}
  \end{center}}

% --- Trang bìa: \trangbia{tiêu đề}{dòng phụ}{phiên bản}
% Chỉ có chữ: tiêu đề, dòng phụ, phiên bản, một dòng ghi chú ở cuối trang.
\newcommand{\trangbia}[3]{%
  \begin{titlepage}
  \vspace*{5.5cm}
  {\sffamily\bfseries\fontsize{40}{44}\selectfont\color{muc}#1\par}
  \vspace{12pt}
  {\sffamily\Large\color{chu-mo}#2\par}
  \vspace{20pt}
  {\sffamily\bfseries\color{sumi}Phiên bản #3\par}
  \vfill
  {\color{vien}\rule{\linewidth}{0.6pt}}\par\vspace{6pt}
  {\sffamily\small\color{chu-mo}Tài liệu chỉ đọc, luôn là bản mới nhất.
    Bài tập nằm trong sân tập riêng của bạn.\par}
  \end{titlepage}}
